^^A Appendix: page breaks around headings, lists and paragraph lines.
^^A To be \input inside didactic.dtx, where % is ignored by \DocInput: this file
^^A therefore has no % comments at all, only ^^A, and contains no literal
^^A vertical bar (the driver makes it a short-verb character).
^^A Needs: biblatex (or plain \cite with a style that handles @online), csquotes,
^^A booktabs, longtable, cleveref.  Assumes \appendix has been issued, or issue
^^A it directly above \input.  didactic provides sidecaption in article.
^^A
^^A Author's note (verification and reproduction):
^^A   scholar session : heading-foot-of-page
^^A   exports         : heading-foot-of-page.csv, appendix-hitlist.tex,
^^A                     queries/q*-*.csv, all in this repository
^^A   searched        : 2026-11-03; providers openalex wos scopus ieee arxiv hal,
^^A                     -l 60 per provider and query; s2 (key rejected, HTTP 403)
^^A                     and dblp (non-JSON, HTTP 202 text/html) failed the health
^^A                     check and were not run
^^A   query numbers   : S1-S6 = q1-q6, S7 = q14, S8 = q15, S9 = q22, S10 = q23,
^^A                     S11 = q12; C1-C3 = q7-q9, C4 = q10, C5 = q11, C6 = q13,
^^A                     C7 = q16, C8-C12 = q17-q21, C13 = q24
^^A   classification  : scholar llm context + llm classify --no-examples --no-enrich,
^^A                     batches of 25, model openai-codex/gpt-5.6-sol, tags
^^A                     supports-claim qualifies-claim adjacent-subtopic
^^A                     off-topic-false-hit; 1238 records.  Model labels for 1171
^^A                     records stand unreviewed.  67 reviewed with scholar sessions
^^A                     decide: all 50 the model kept, all 15 discards under 0.7,
^^A                     and the 2 general reviews cited.  Abstracts present for
^^A                     1044 of 1238 records; enrichment did not finish (s2 dead,
^^A                     openalex timeouts), so the rest went by title.
^^A                     session-before-review.json is the state before the review.
^^A   hit-list table  : scholar sessions export heading-foot-of-page -f table
^^A                     --bearing-only --lang en --label pb-hits
^^A                     --track "Page breaks around headings, lists and paragraph
^^A                     lines" --theme lines-after-heading="lines after a heading"
^^A                     --theme space-at-foot="space at the page foot"
^^A                     --theme reader-evidence="reader evidence"
^^A                     -o appendix-hitlist; then sed -i '/^%/d' on the result,
^^A                     since \DocInput would typeset its % comment header.
^^A   provenance      : didactic.bib, every entry
^^A   mechanism test  : lit/test/t-pen.py (pdflatex, article, TeX Live 2025)
^^A   chicago manual  : read online through KTH (SOCKS tunnel, curl with a cookie
^^A                     jar, no credentials); see closed-access.md in backing-claims
^^A   SECONDARY       : older Chicago editions via mittelbach2018forlorn;
^^A                     Bringhurst via izma2023widows;
^^A                     Tschichold and Wheildon via wilson2018notes; Wheildon via
^^A                     harrison2020widows
^^A
^^A Wiring: decisions go in their own section, not here.  It should point to
^^A app:pb-lines, app:pb-gap, app:pb-lists and app:pb-penalties in a footnote at
^^A each claim, saying what the appendix answers.

\providecommand{\apcs}[1]{\texttt{\textbackslash #1}}

\section{Page breaks: questions, answers and method}\label{app:pb-method}

This appendix reports what the typographic and \LaTeX\ literature says about
four questions on page breaks.
They arise for any package that decides where a page ends in lecture notes.
The appendix gives the evidence only; it does not recommend a design.
What this package decided from it, and why, is in \cref{PageBreakDecisions}.
\begin{description}
\item[\Cref{app:pb-lines}] How many lines must follow a heading at the foot of
  a page?  Where a number is stated, at least two, including in the University of
  Chicago Press's sample house style; three occurs only as the top of a range,
  \enquote{two or three}.
\item[\Cref{app:pb-gap}] When they do not fit, does the heading move, and how
  much empty space is accepted for that?  Where sources say, it moves.  One source
  sets a limit for that case: up to four lines short, in the same sample house
  style.
\item[\Cref{app:pb-lists}] Should a list's lead-in stay with the list, and how
  are lists treated at a page break?  The same sample asks for two lines on each side of a
  break, in lines and not items.  Other evidence is thin and points both ways.
\item[\Cref{app:pb-penalties}] How strongly should a package forbid a first or
  last line alone at a page break?  No source compares the values for readers.
  They differ mainly in what they do when a page has no flexibility.
\end{description}

\paragraph{What was not read.}
Of the classical handbooks, only \emph{The Chicago Manual of Style} could be read,
in three editions.\cite{chicago2010manual,chicago2017manual,chicago2024manual}
The others could not be opened at the relevant passages: \emph{New Hart's Rules},
Butcher's \emph{Copy-editing}, Bringhurst's \emph{The Elements of Typographic Style},
Tschichold's \emph{The Form of the Book}, McLean's \emph{Thames \& Hudson Manual
of Typography}, Lamport's \emph{\LaTeX: A Document Preparation System} and
\emph{The \LaTeX\ Companion}.
The answers below therefore rest on one handbook and on other kinds of source, and
any statement about the unread books is second hand and marked as such.
Checking them is the first item of future work.

\paragraph{Databases.}
The searches were run on 3~November 2026 with the command-line tool
\texttt{scholar}\cite{bosk2026scholar}, on six databases that passed its health
check that day: OpenAlex, Web of Science, Scopus, IEEE Xplore, arXiv and HAL.
Two could not be searched: Semantic Scholar rejected the API key (HTTP 403), and
DBLP answered with a web page instead of data (HTTP 202).
Every query ran on all six databases, with the same keywords, translated into
each database's field syntax where it needs one (\cref{tab:pb-queries}).
OpenAlex ranks by relevance and returned far more noise than the others, and
counts are capped at 60 per database and query.

^^A Longtable, not the guide's sidecaption float: 24 wrapped rows do not fit one
^^A page of the article layout, and the fallback sidecaption's caption was lost.
\begingroup\scriptsize
\setlength{\tabcolsep}{2.5pt}
\begin{longtable}{@{}l>{\raggedright\arraybackslash}p{0.43\linewidth}rrrrrrrr@{}}
\caption{Queries and hits per database, 3~November 2026, tool
    \texttt{scholar}.  S = support query, C = counter query.  Databases:
    OpenAlex (OA; not open access), Web of Science (WoS), Scopus, IEEE Xplore
    (IEEE), arXiv, HAL, Semantic Scholar (S2), DBLP.  At most 60 hits per
    database and query, so 60 means \enquote{at least 60}.  \enquote{--} means
    the search failed: S2 rejected the key and DBLP answered with a web page, so
    neither could be searched.  Web of Science received the keywords as
    \texttt{TS=(\dots)}, Scopus as \texttt{TITLE-ABS-KEY(\dots)}, the others as
    boolean free text.}\label{tab:pb-queries}\\
\toprule
& Keywords & OA & WoS & Scopus & IEEE & arXiv & HAL & S2 & DBLP \\
\midrule
\endfirsthead
\multicolumn{10}{@{}l}{\emph{\tablename~\thetable{} (cont.)}}\\
\toprule
& Keywords & OA & WoS & Scopus & IEEE & arXiv & HAL & S2 & DBLP \\
\midrule
\endhead
\bottomrule
\endlastfoot
      S1 & \texttt{"widows and orphans" AND (typography OR typesetting OR "page layout")} & 57 & 1 & 1 & 0 & 1 & 0 & -- & -- \\
      S2 & \texttt{heading AND "page break" AND (typography OR typesetting OR pagination)} & 60 & 0 & 0 & 0 & 0 & 0 & -- & -- \\
      S3 & \texttt{pagination AND typesetting AND ("white space" OR "short page" OR "facing pages")} & 60 & 0 & 0 & 1 & 0 & 0 & -- & -- \\
      S4 & \texttt{list AND "page break" AND (typography OR typesetting OR pagination)} & 60 & 0 & 1 & 0 & 0 & 0 & -- & -- \\
      S5 & \texttt{"page breaking" AND (TeX OR LaTeX)} & 60 & 1 & 3 & 0 & 1 & 0 & -- & -- \\
      S6 & \texttt{("widow line" OR "orphan line" OR "club line" OR Hurenkind OR Schusterjunge)} & 60 & 0 & 3 & 0 & 1 & 0 & -- & -- \\
      S7 & \texttt{"short page" OR "ragged bottom" OR "running short" OR "flush bottom" OR "bottom balancing"} & 60 & 60 & 60 & 39 & 12 & 7 & -- & -- \\
      S8 & \texttt{"memoir" AND LaTeX AND (class OR typesetting OR typography)} & 60 & 0 & 0 & 0 & 1 & 0 & -- & -- \\
      S9 & \texttt{(subheading OR subhead OR "section heading") AND ("page break" OR pagination OR "keep with next" OR "foot of the page" OR "bottom of the page") AND (typography OR typesetting)} & 59 & 0 & 0 & 0 & 0 & 0 & -- & -- \\
      S10 & \texttt{("widow penalty" OR "orphan penalty" OR "club penalty" OR clubpenalty OR widowpenalty OR "widow control" OR "orphan control") AND (typesetting OR pagination OR TeX OR LaTeX)} & 39 & 0 & 0 & 0 & 0 & 0 & -- & -- \\
      S11 & \texttt{("bulleted list" OR "list items" OR "lead-in") AND ("page break" OR pagination) AND (reading OR comprehension OR typography)} & 60 & 0 & 0 & 0 & 2 & 0 & -- & -- \\
      \midrule
      C1 & \texttt{("widows and orphans" OR "page break") AND (reading OR comprehension OR legibility OR readability) AND (experiment OR empirical OR evidence)} & 60 & 0 & 1 & 0 & 0 & 0 & -- & -- \\
      C2 & \texttt{typography AND (rules OR conventions) AND (empirical OR evidence OR myth OR critique) AND (readability OR legibility)} & 60 & 1 & 0 & 0 & 0 & 0 & -- & -- \\
      C3 & \texttt{("page turn" OR "page break" OR "page boundary") AND ("eye movements" OR comprehension)} & 60 & 2 & 1 & 0 & 0 & 0 & -- & -- \\
      C4 & \texttt{(heading OR "widows and orphans" OR "page break") AND typesetting AND (criticism OR overrated OR "trade-off" OR "white space")} & 60 & 0 & 0 & 1 & 1 & 0 & -- & -- \\
      C5 & \texttt{("page break" OR "page turn" OR "page boundary") AND (sentence OR paragraph) AND (reading OR memory OR recall)} & 60 & 1 & 1 & 0 & 0 & 0 & -- & -- \\
      C6 & \texttt{(heading OR headings) AND ("page layout" OR "white space" OR "page break") AND (reading OR comprehension) AND experiment} & 60 & 0 & 0 & 0 & 0 & 0 & -- & -- \\
      C7 & \texttt{"memoir" AND LaTeX AND (criticism OR limitations OR "page break" OR "ragged bottom" OR "flush bottom")} & 59 & 0 & 0 & 0 & 2 & 0 & -- & -- \\
      C8 & \texttt{("widows and orphans" OR "widow line" OR "orphan line" OR "club line" OR Hurenkind OR Schusterjunge) AND (criticism OR empirical OR evidence OR experiment OR reading OR comprehension)} & 60 & 7 & 8 & 0 & 1 & 1 & -- & -- \\
      C9 & \texttt{("short page" OR "ragged bottom" OR "running short" OR "flush bottom" OR "bottom balancing" OR "facing pages") AND (typography OR typesetting OR pagination) AND (criticism OR empirical OR experiment OR reading OR comprehension)} & 60 & 0 & 1 & 0 & 0 & 0 & -- & -- \\
      C10 & \texttt{(list OR "bulleted list" OR "list items" OR "lead-in") AND ("page break" OR pagination) AND (criticism OR empirical OR experiment OR evidence)} & 60 & 2 & 4 & 2 & 6 & 0 & -- & -- \\
      C11 & \texttt{("page breaking" OR pagination) AND (TeX OR LaTeX) AND (criticism OR limitations OR empirical OR evaluation OR "white space")} & 60 & 1 & 3 & 0 & 2 & 0 & -- & -- \\
      C12 & \texttt{(heading OR subheading OR subhead OR "section heading") AND ("page break" OR pagination OR "keep with next" OR "foot of the page" OR "bottom of the page") AND (reading OR comprehension OR experiment OR empirical OR criticism)} & 60 & 0 & 2 & 0 & 0 & 0 & -- & -- \\
      C13 & \texttt{("widow penalty" OR "orphan penalty" OR "club penalty" OR clubpenalty OR widowpenalty OR "widow control" OR "orphan control") AND (criticism OR empirical OR experiment OR reading OR "white space")} & 60 & 0 & 1 & 0 & 0 & 0 & -- & -- \\
\end{longtable}
\endgroup

\paragraph{Web searches.}
This literature is mostly books, manuals and community documents, which the
databases index poorly, so web searches complemented them.
\Cref{tab:pb-web} lists them verbatim.
Each source they led to was then opened; a search snippet was never used as
evidence.

\begingroup\scriptsize
\begin{longtable}{@{}>{\raggedright\arraybackslash}p{0.62\linewidth}>{\raggedright\arraybackslash}p{0.34\linewidth}@{}}
\caption{Web searches, 3~November 2026, with the outcome of each.}\label{tab:pb-web}\\
\toprule
Query & Outcome \\
\midrule
\endfirsthead
\multicolumn{2}{@{}l}{\emph{\tablename~\thetable{} (cont.)}}\\
\toprule
Query & Outcome \\
\midrule
\endhead
\bottomrule
\endlastfoot
      \texttt{Chicago Manual of Style page makeup "subhead" bottom of page "two lines of text" must follow widow first line of paragraph} & No Manual text; pointers only \\
      \texttt{heading at bottom of page how many lines of text must follow typography rule at least two lines after subheading book typesetting} & A practitioner thread\cite{binder2023lonely} \\
      \texttt{widows and orphans evidence reading research do they matter overrated typography no empirical study} & A secondary report\cite{harrison2020widows}; no study \\
      \texttt{"subhead" "at least two lines" foot of page Chicago Manual of Style "2." page makeup "must be followed by"} & Publisher's blog\cite{cmosshoptalk2018subheads} \\
      \texttt{New Hart's Rules Butcher's Copy-editing heading bottom of page "at least two lines" text below heading page make-up short page} & Nothing usable \\
      \texttt{Überschrift am Seitenende mindestens drei Zeilen Text folgen Typografie Regel Mikrotypografie Satzregeln Schusterjunge} & No three-line rule verified \\
      \texttt{rubrik längst ner på sidan minst två rader text efter rubriken typografi regel sidbrytning svenska skrivregler horunge} & No authority found \\
      \texttt{titre en bas de page doit être suivi d'au moins deux lignes de texte règles typographiques mise en page lexique imprimerie nationale} & No number verified \\
      \texttt{bulleted list page break lead-in sentence keep with first item minimum items on a page technical writing style guide} & A practitioner thread\cite{sims2021listtags} \\
      \texttt{"A page should not begin with the last line of a paragraph" subhead "two lines" page makeup Chicago Manual} & The W3C draft\cite{cramer2016latinreq} \\
      \texttt{Bringhurst Elements of Typographic Style "Balance facing pages by moving single lines" widows "Leave at least two lines" subhead} & Second-hand quotations only \\
      \texttt{Turabian thesis formatting subhead bottom of page "at least two lines of text" below subheading must follow} & College guides\cite{georgefox2026turabian,jsums2024turabian,nobts2026template} \\
      \texttt{tex.stackexchange section heading at bottom of page keep with at least three lines needspace titlesec nobottomtitles how many lines after heading} & Forum uses, not advice \\
      \texttt{page bottom white space short page "facing pages" one line short heading moved next page typesetter accept book design rule running pages short} & A vendor guide\cite{cambric2026widow} \\
      \texttt{InDesign Keep Options "Keep With Next" lines "up to five" headings paragraph Adobe help} & Adobe's help\cite{adobe2026incopy} \\
      \texttt{TeX FAQ widows and orphans clubpenalty "section" heading page break "at bottom of page" texfaq; LaTeX Companion needspace "subheading" page break recommendation Mittelbach} & The \TeX\ FAQ\cite{texfaq2026widows}; a book author's notes\cite{osborne2010pmakeup} \\
      \texttt{study reader effect sentence split across page turn reading comprehension eye tracking page boundary paragraph interrupted text print} & One screen study\cite{dillon1990splitting} \\
      \texttt{"Towards evidence-based typography" literature review experiment design widows orphans typographic rules not empirically tested} & Two general reviews\cite{veytsman2011evidence,vanderwaarde2025nineteen} \\
      \texttt{thesis dissertation formatting manual heading bottom of page "at least two lines" OR "at least three lines" OR "two lines of text" of text following heading university graduate school} & Thesis manuals\cite{pitt2026etd,utah2026text,iowa2020manuscript,uwest2026format}; no three-line rule \\
      \texttt{memoir class \textbackslash raggedbottomsection \textbackslash bottomsectionskip short page section heading moved flushbottom stretched} & Copies of the manual only \\
      \texttt{memoir LaTeX class default \textbackslash clubpenalty 1000 \textbackslash widowpenalty 1000 why versus 150 widows orphans still appear criticism} & No critique of the defaults \\
\end{longtable}
\endgroup

\paragraph{Known documents.}
Specifications, package manuals and program sources were retrieved directly,
which is the right method for them: they are not journal articles and have no
place in these databases.
They were read in full at the cited passages.
Package manuals and sources are the ones in \TeX\ Live 2025.
\emph{The \TeX book}\cite{knuth1986texbook} was read in a copy in the maintainer's
own library, which was also searched by title for the handbooks above without
result.
\emph{The Chicago Manual of Style} was read in its 18th and 17th editions online,
through a university's institutional access, at the sections named by its index, and
in its 16th edition in a PDF in the maintainer's own files.
A text search of a PDF does not see figures, so the page-makeup figure of the 16th
edition was read as an image, from a pointer in its text.
\emph{New Hart's Rules}, Butcher and \emph{The \LaTeX\ Companion} were tried through
the same access and refused or offered for purchase; search inside the Internet
Archive's lending copies was also refused.

\paragraph{Screening.}
All 1238 records were labelled by a language model against a statement of the four questions, into four
categories: supports, qualifies, adjacent and off-topic.
These labels are the model's and were not all checked.
The model worked from title and abstract; 194 records had no abstract and went by title alone.
Sixty-seven records were reviewed by hand, from title and abstract: all 50 the model kept, all 15 discards below 0.7
confidence, and two general reviews it had called adjacent.
The other 1171 discards stand as the model left them, apart from a keyword check of their titles and abstracts for page-break
vocabulary, which changed nothing.
The review moved 26 of the model's 50 kept records to adjacent or off-topic, among them all three it had called
qualifying: a study of picture books, a study of test items and a note on old book pagination.
\Cref{tab:pb-hits} at the end of this appendix lists the records that bear on the questions.
A row marked \enquote{supports} there was judged relevant from its title or abstract and, unless cited, was not read: it is a
pointer for further reading and not evidence used here.
Most sources cited in this appendix are known documents and web results, which are not in that table.
The provenance of each source is recorded in the bibliography's source file with
fixed fields: the claim, how it was found, why it was picked, a verbatim quote,
what was read, the counter-search and the date.

\paragraph{Mechanism test.}
Where sources describe what \TeX\ does, the behaviour was also checked with small
documents in the \texttt{article} class, on a page of 20 lines, under
\texttt{pdflatex} in \TeX\ Live 2025.
This checks the mechanism only and says nothing about readers.

\paragraph{Shared limitations.}
Semantic Scholar and DBLP were not searched.
A language model's labels decided most of the screening, so a relevant record may have been discarded unseen.
OpenAlex counts are capped and ranked, so relevant records may lie below the cap.
Books and pre-digital reading research are poorly indexed, so an empty database
result is weak evidence of absence.
Web results depend on the search engine and were not exhaustive.
Thesis manuals, package manuals and practitioner posts state what their authors
require or prefer, not what benefits readers.

\section{Lines after a heading}\label{app:pb-lines}

\paragraph{Claim and question.}
A heading should not stand alone at the foot of a page; the W3C requirements for
Latin text layout put this as a minimum number of following
lines.\cite{cramer2016latinreq}
The question is: how many lines of text must follow a heading at the foot of a
page?
The strength tested is a stated minimum, as distinct from a general wish to keep
heading and text together.

\paragraph{Method.}
See \cref{app:pb-method}.
Queries S1, S2, S6 and S9 carry the support side, C1, C4, C6, C8 and C12 the
counter side.

\paragraph{Results: stated numbers.}
The strongest source is a publisher's own rule.
The 16th edition of \emph{The Chicago Manual of Style} prints, as a sample of the rules a
publisher gives typesetters,\cite[§A.17]{chicago2010manual} a \enquote{House Style for Composition and
Page Makeup} of the University of Chicago Press: \enquote{The minimum number of lines of text below a subhead
at the bottom of a page is two.}\cite[fig.~A.3, Page Makeup, item 4]{chicago2010manual}
The only source seen that states a range is the W3C draft, edited by a publisher's
employee and based on print practice: \enquote{A head should never be the last
thing on a page; it must be followed by two or three lines of
text.}\cite[§10.3]{cramer2016latinreq}
The same sentence stands in the published Working Draft.\cite[§9.3]{cramer2014latinreq}
It gives no source for the number.
Four university thesis manuals set two.
Pittsburgh: \enquote{A subheading at the bottom of a page should be followed by at
least two full lines of type.}\cite[Margins]{pitt2026etd}
Iowa: \enquote{Headings may not be placed at the bottom of the page without at
least two lines of text below them.}\cite{iowa2020manuscript}
Utah and the University of the West say the
same.\cite{utah2026text,utah2026checklist,uwest2026format}
Theses are double or one-and-a-half spaced, so these two lines are taller than two
lines of a book.

\paragraph{Results: no number.}
The running text of the Manual gives no count.
The 18th edition says: \enquote{A page should not end with a subhead}, and goes on: \enquote{A page can,
however, end with the first line of a new paragraph}.\cite[§2.121]{chicago2024manual}
The 17th edition has the same sentences.\cite[§2.116]{chicago2017manual}
The 16th has the second sentence without the first.\cite[§2.113]{chicago2010manual}
Taken alone, the current text would not forbid a subhead followed by a single line, but that is this
review's inference, and the Press's sample sets two.
No such figure was found through the 18th edition's index, so whether the two-line rule still stands there is
not known.
The publisher's blog says, for student papers: \enquote{Never end a page with a
subhead.}\cite{cmosshoptalk2018subheads}
College guides repeat this.\cite{georgefox2026turabian,jsums2024turabian}
One seminary's template counts a subheading alone at the page bottom as an orphan
error.\cite{nobts2026template}
A guide for scientific texts asks that a heading style make sure a heading
\enquote{is not followed by a page break}.\cite{schwen2020tensimple}
Butterick advises keep-with-next for headings.\cite{butterick2026headings}
A TUGboat article calls separating a section heading from its paragraph \enquote{almost
always undesirable}.\cite[§9.6]{chernoff2022lwc}
A W3C statement of the reason is the only one found: \enquote{A heading at the bottom
of a page removes the title from the section, and the section from the
title.}\cite[§2]{cramer2016latinreq}
None of these is backed by reader data.

\paragraph{Results: systems, as evidence of practice.}
\LaTeX\ forbids a break after the first line of the paragraph following a heading,
through \apcs{clubpenalty} set to 10000 in \apcs{@afterheading}.\cite{latex2025sources}
The mechanism test agrees: with a paragraph of several lines, heading and two lines
stayed at the foot, and with room for only one line the heading moved.
A one-line paragraph gets no such protection: in the test, a heading followed by a
one-line paragraph stayed at the foot, with the next paragraph on the next page.
The memoir class keeps this mechanism, with the same results in the test, and its manual gives no
number.\cite{wilson2025memoir,wilson2025memoirsrc}
CSS's sample style sheet avoids a break after headings, and its \texttt{orphans}
property defaults to 2, which together imply two
lines.\cite[Appendix~D, §13.3.2]{bos2011css21}\cite[§3.3]{atanassov2018cssbreak}
XSL-FO has the same default, and the DocBook style sheets keep section titles with what
follows.\cite[§7.20.6]{berglund2006xsl11}\cite{docbook2026sectiontitle}
These two-line figures are this review's inference from combining rules, not statements
in the specifications.
Adobe's software lets the user set up to five lines and describes the purpose as keeping
a heading with \enquote{the next few lines}.\cite{adobe2026incopy}
Word's keep-with-next states no count.\cite{microsoft2026keep}
The \texttt{nowidow} package defaults to two, for paragraphs.\cite[§2.2]{pinson2011nowidow}
The \texttt{needspace} package leaves the amount to the user.\cite{wilson2025needspace}

\paragraph{Results: other positions.}
On the low side, an Adobe community expert calls two \enquote{pretty standard}, adding
\enquote{I do have one client who prefers the single line}, and later \enquote{choosing
between 1 or 2 lines is just a preference}.\cite{binder2023lonely}
On the high side, one user of another system wanted a heading with \enquote{three lines of
text} at the page bottom pushed to the next page.\cite{bryon2026widowing}
Both are anecdotes.
Searches in German, Swedish and French verified no rule of their own.

\paragraph{Discussion and limitations.}
The figure in the Manual is labelled a sample and is the Press's own house style, not a rule the Manual
lays down for all publishers; it nevertheless shows what one major book publisher requires.
The other sources that state a number are a draft requirements document and thesis manuals.
Chicago's running text states no number, which is silence, not a lower floor.
The other handbooks were not read and may say something else.
The W3C draft is the only source for three, and only as the top of a range.
Agreement among system defaults is weak evidence, since they may share an origin.
No study was found on how many lines readers need under a heading.

\paragraph{Conclusion.}
Where a number is stated, the floor in the literature read is two lines; three appears
only in \enquote{two or three}, and no source read requires three.
The University of Chicago Press's sample house style states exactly two, and the Manual's running text
states no number.
The support for two is moderate to good for a craft rule and the support for three is weak.
That a heading must not stand alone has broad support, but as craft practice and not as a
measured effect.
What this package does with it: \cref{PageBreakDecisions}.

\paragraph{Future work.}
This search left \emph{New Hart's Rules}, Butcher's \emph{Copy-editing} and Bringhurst unread at
this point, did not establish whether the 17th and 18th editions of the Manual keep the figure, and it left Semantic Scholar and DBLP unsearched.
The field lacks any experiment on the number of lines under a heading: a study varying zero, one,
two and three lines at the foot, measuring rereading and navigation, would settle whether the rule
has a reader effect at all.
If one line proved harmless, the two-line floor would stand as convention only.

\paragraph{Hit list.}
\Cref{tab:pb-hits} lists the database records; of those, Schwen and Chernoff are cited here.

\section{The trade-off with empty space}\label{app:pb-gap}

\paragraph{Claim and question.}
A heading that lacks room for its lines moves to the next page, at the cost of a short
page.\cite{utah2026text}
The question is: when the lines do not fit, does the heading move, and how much empty space does
the literature accept for that?

\paragraph{Method.}
See \cref{app:pb-method}.
Queries S3, S7 and S8 carry the support side, C4, C7 and C9 the counter side.

\paragraph{Results: the heading moves.}
Every source that says what happens to the heading moves it.
Utah: \enquote{If there is not enough room, the subheading moves to the top of the next page.
This is the only time a gap is allowed in the text when it is not the end of a chapter or before
an equation too long to fit at the bottom of a page.}\cite[Placement]{utah2026text}
Its checklist also allows a footnote-forced gap, so \enquote{only} is not
strict.\cite[Margins]{utah2026checklist}
Pittsburgh: \enquote{the subheading should begin the next page.}\cite[Margins]{pitt2026etd}
Memoir describes the same behaviour of the standard classes: a subhead \enquote{too close to the
bottom of a page is moved to the top of the following page}.\cite[§6.2, p.~75]{wilson2025memoir}
Chernoff's package does so by default, with this cost: \enquote{moving the section heading can create
a large blank space at the end of the page.}\cite[§9.6, p.~34]{chernoff2022lwc}

\paragraph{Results: how much space.}
One source sets a limit for exactly this case.
The University of Chicago Press's sample: \enquote{The text page may be up to four lines short if there is no room
for a subhead at the bottom of a page. The subhead begins at the top of the next page.}\cite[fig.~A.3, item 4]{chicago2010manual}
It also says that chapter-opening pages \enquote{may be up to two lines short} to avoid a bad break.\cite[fig.~A.3, item 3]{chicago2010manual}
The same figure limits other adjustments more tightly: \enquote{Facing text pages should be of equal depth and must align}, and
spreads end \enquote{either at normal depth, or exactly one line short or one line long}.\cite[fig.~A.3, items 6, 8]{chicago2010manual}
It does not say how four lines fit with this.
The Manual's running text gives the one-line figure too: \enquote{it is usually acceptable for both pages to run a line
long or short to avoid widows}, or to accommodate corrections,\cite[§2.122]{chicago2024manual} and for a subhead at the foot and the
like, \enquote{page length may be adjusted}.\cite[§2.121]{chicago2024manual}
The other numbers measure other things.
The \texttt{titlesec} option \texttt{nobottomtitles} moves titles when little room is left, with a
threshold \enquote{whose default value is .2\apcs{textheight}}, but it is off by default and is a
package setting, not advice.\cite[p.~6]{bezos2025titlesec}
Plain \TeX's \apcs{beginsection} places \texttt{plus.3\apcs{vsize}} of stretch before a section;
reading this as a tolerance is this review's interpretation.\cite{knuth2021plain}
A vendor guide calls a page \enquote{three or four lines short} arguably worse than the widow it
fixes.\cite{cambric2026widow}
One book author reports publisher rules allowing pages \enquote{longer or shorter than x by up to one
line}.\cite{osborne2010pmakeup}
Wilson's notes give \enquote{Five lines on the last page of a chapter} as a reasonable minimum, which
concerns chapter ends.\cite[§3.3.3, p.~52]{wilson2018notes}

\paragraph{Results: short pages as a defect.}
Book-oriented sources treat the short page as a real cost.
Mittelbach: \enquote{all \LaTeX\ can do is run the page or column short, generating a fairly ugly hole
at the bottom.}\cite[p.~247]{mittelbach2018forlorn}
The \TeX\ FAQ: \enquote{many publishers insist on the default \apcs{flushbottom}; it is seldom acceptable
to introduce stretchability into the vertical list, except at points (such as section headings) where
the document design explicitly permits it.}\cite{texfaq2026widows}
Wilson accepts space around headings, \enquote{a heading is meant to attract attention}, but calls too
much white space in the typeblock \enquote{a distraction to the reader}.\cite[§3.3.1, p.~48]{wilson2018notes}
The Press's sample forbids one common fix: \enquote{Adding extra space between lines or paragraphs to make facing pages align
is not allowed}, though it allows small adjustments around separate text elements.\cite[fig.~A.3, Typography item 5, Page Makeup item 8]{chicago2010manual}
Memoir warns that a flush bottom \enquote{may have a deleterious effect on the color of some pages} and
offers \apcs{raggedbottomsection}, off by default, for pages short because of a moved
subhead.\cite[pp.~6, 76]{wilson2025memoir}
It neither recommends nor discourages the switch.
Butterick finds blank lines \enquote{more annoying than the widows and
orphans},\cite{butterick2026widow} but also keeps heading lines together, warning of
gaps.\cite{butterick2026keeplines}

\paragraph{Results: facing pages, and book versus report.}
The main distinction is facing pages, as in Chicago above.
W3C: \enquote{If a document has facing pages, the publisher usually requires that they align top and
bottom}, with exceptions including \enquote{When aligning facing pages will make some other issue
worse}; running short means removing a line from each page of the
spread.\cite[§§5.3--5.4]{cramer2016latinreq}
KOMA-Script: \enquote{In two-sided documents that is usually undesirable.}\cite[§2.1, p.~27]{kohm2025koma}
Osborne: \enquote{Every pair of facing pages (even--odd) must have the same length.}\cite{osborne2010pmakeup}
The standard classes and memoir encode the split, with a ragged bottom only when
one-sided.\cite{latex2025classes}\cite[p.~6]{wilson2025memoir}
Book versus report is only implicit: thesis manuals accept the gap and trade-book sources resist it, but
no source read draws the contrast.

\paragraph{Results: fewer lines to avoid space.}
No source read argues, as a principle, that a heading should keep fewer lines to avoid white space.
The nearest are the Adobe expert's \enquote{just a preference}\cite{binder2023lonely} and a
\texttt{lua-widow-control} option that leaves the orphan under the heading and warns.\cite[§9.6]{chernoff2022lwc}
The U.S. Government Publishing Office makes the analogous trade for paragraphs only: widows
\enquote{are permitted if absolutely necessary to maintain uniform makeup and page
depth}.\cite[§2.6]{gpo2016style}
It has no rule on heads at the foot that a keyword search could find.

\paragraph{Discussion and limitations.}
Thesis manuals concern double-spaced work with fixed margins, and say nothing about how a book looks.
The numbers above measure different things and should not be compared.
Printed spreads and one-sided or on-screen reading may differ, and the sources do not address screens.
No reader study of short pages was found.

\paragraph{Conclusion.}
Yes, the heading moves in every source that says what happens to it, and the only stated limit for that case is
\enquote{up to four lines short}, in the University of Chicago Press's sample house style.
For other adjustments the same Press allows one line, on both facing pages.
Short pages are treated as a defect in book work, mainly where facing pages must align, and as a
permitted exception in thesis manuals, which set no limit.
The strength is moderate for the move and limited for the amount, which rests on one publisher's sample.
What this package does with it: \cref{PageBreakDecisions}.

\paragraph{Future work.}
Besides the unread handbooks of \cref{app:pb-method}, Bringhurst's rule on balancing facing pages should be read
at source; a W3C appendix lists \enquote{Balance facing pages by moving single lines} without
attribution.\cite{cramer2016latinreq}
Tschichold's reported view, that a page a line shorter or longer is less of an affront than a widow, is
known here only through Wilson.\cite[§3.3.3]{wilson2018notes}
The field lacks any measure of whether readers notice or mind a short page, on spreads or single pages;
a study would show whether the concern is typographic only.

\paragraph{Hit list.}
\Cref{tab:pb-hits} lists the database records; of those, Chernoff is cited here.

\section{Lists at a page break}\label{app:pb-lists}

\paragraph{Claim and question.}
A break between a lead-in sentence and its list is said to be annoying.\cite[§18.7, p.~333]{wilson2025memoir}
The question is: should a list's lead-in stay with the list, and how are lists treated at a page break?

\paragraph{Method.}
See \cref{app:pb-method}.
Queries S4 and S11 carry the support side, C10 the counter side.

\paragraph{Results: the publisher's sample.}
The University of Chicago Press's sample house style is the only book-publishing source found with a rule, under
\enquote{Line minimums for breaking elements}: \enquote{A minimum of two text lines must appear above or below a separate text element such as
an extract, poem, or list}, and \enquote{A minimum of two lines of a separate text element may fall at the top or bottom of a
page.}\cite[fig.~A.3, item 5]{chicago2010manual}
So a list may break across pages, and the count is in lines, not items.
It does not name the lead-in; that two text lines above would normally include it is this review's reading.

\paragraph{Results: the lead-in.}
Memoir: \enquote{A sentence or two may be used to introduce a list (e.g., itemize) and it can be annoying
if there is a page break between the introductory words and the first item.}\cite[§18.7, p.~333]{wilson2025memoir}
It offers \apcs{noprelistbreak}, per list.
CSS's sample style sheet has \texttt{ul, ol, dl \{ page-break-before: avoid \}}.\cite[Appendix~D]{bos2011css21}
A FrameMaker user writes that \enquote{the first list item has to stay with the preceding paragraph because
it's an introduction into the list}.\cite{sims2021listtags}
Mittelbach reports, for a nearby case, that a display formula starting a page with its introduction on the
previous page \enquote{is considered a no-go by nearly everybody}.\cite[p.~247]{mittelbach2018forlorn}
This is an analogy, not a rule about lists.

\paragraph{Results: the systems that do the opposite.}
The standard \LaTeX\ classes set \apcs{@beginparpenalty} and \apcs{@itempenalty} to
$-51$.\cite{latex2025classes}
Memoir does the same, and its source says why: \enquote{They are set to a bonus value to encourage page breaking
at these points.}\cite{wilson2025memoirsrc}
So the default encourages a break at exactly the place memoir's manual calls annoying.
A list directly after a heading is exempt.
DocBook's list properties set no keep.\cite{docbook2026sectiontitle}

\paragraph{Results: items at a break.}
One practitioner reports: \enquote{Our style guide mandates at least 2 list items at the top or bottom of a
page if the list can't be kept all on one page. The only exception is if an item is long enough to split over
the page break.}\cite{sims2021listtags}
That is one unnamed in-house guide, second hand, and it counts items where the Press's sample counts lines.
No source read treats items split across pages as a separate topic.

\paragraph{Results: the handbook.}
\emph{The Chicago Manual of Style} treats lists in its sections on punctuation and
format.\cite[§§6.138, 6.141--6.142]{chicago2024manual}
A keyword search of those sections found no rule on page breaks, its index has no such entry under lists, and
the bad breaks it lists for proofreading do not include lists.\cite[§2.121]{chicago2024manual}

\paragraph{Discussion and limitations.}
Apart from the Press's sample the evidence is thin: one manual's remark, one informative style sheet, forum posts.
The sample's wording leaves open how \enquote{above or below} applies when a list starts a page.
The silence of Chicago's running text is a keyword search of the sections named, not a reading of the whole book.
Keyword searches of the package manuals found nothing further.
No reader study was found.

\paragraph{Conclusion.}
One publisher's sample requires two text lines beside a list and two lines of a list on each side of a break; keeping the lead-in itself
with the list has only weak, indirect support, and the standard \LaTeX\ defaults point the other way.
A minimum counted in items has one second-hand report of two.
None of this is backed by reader data, so the evidence does not settle the question.
What this package does with it: \cref{PageBreakDecisions}.

\paragraph{Future work.}
The unread handbooks of \cref{app:pb-method} and technical-writing manuals such as the \emph{Microsoft Manual of
Style} should be checked for lists at a page break; none was read.
A mailing-list guideline asking for three items on each side of a break could not be opened and is not
counted.
The field lacks any study of whether a lead-in separated from its list, or a lone item, affects reading.

\paragraph{Hit list.}
\Cref{tab:pb-hits} lists the database records; none of those cited in this appendix bears on lists.

\section{How strongly to forbid a lone first or last line}\label{app:pb-penalties}

\paragraph{Claim and question.}
\TeX\ discourages a lone first or last line through penalties.\cite[p.~104]{knuth1986texbook}
The defaults differ: 150 in plain \TeX\ and the \LaTeX\ kernel, 1000 in memoir, and 10000, which forbids the
break, in many journal classes.\cite{knuth2021plain,latex2025sources,wilson2025memoirsrc}\cite[p.~246]{mittelbach2018forlorn}
The question is: how strongly should a package forbid a first or last line alone at a page break, as the kernel's
150, memoir's 1000, or 10000?

\paragraph{Method.}
See \cref{app:pb-method}.
Queries S1, S5, S6 and S10 carry the support side, C1, C2, C8, C11 and C13 the counter side.

\paragraph{Results: what the values do.}
Knuth's text: such lines \enquote{will not appear all alone on a page unless the alternative is
worse}.\cite[p.~104]{knuth1986texbook}
The cost of a break is badness plus penalty, with badness at most 10000 unless the page is
overfull.\cite[p.~111]{knuth1986texbook}
The \TeX\ FAQ draws the practical consequence: \enquote{for typical layouts, there are only two sensible settings
for the penalties: finite (150 or 500, it doesn't matter which) to allow widows and orphans, and infinite (10000 or
greater) to forbid them.}\cite{texfaq2026widows}
The mechanism test supports this, and bounds it.
On a text-only page of 20 lines with no stretch, 150, 1000 and 9999 all left the orphan in place, and only 10000
moved it, leaving the page a line short.
Memoir with its default of 1000 likewise left the orphan in place.
With a little stretch on the page the values began to differ: with 0.5\,pt of stretch between paragraphs, 150 kept the orphan
and 1000 moved it; with 1\,pt, all three moved it; with 0.2--0.3\,pt, 150 and 1000 kept it and 9999 moved it.
So 150 and 1000 differ only in a narrow band of available stretch.

\paragraph{Results: what the sources recommend.}
For 10000: \enquote{many journal classes for \LaTeX\ completely forbid widows and orphans by setting
\apcs{widowpenalty} and \apcs{clubpenalty} to 10000}, TUGboat included.\cite[p.~246]{mittelbach2018forlorn}
Against: \enquote{The downside, as discussed previously, is partly empty pages and if space is a premium [\dots]
then this is a possible problem. Thus you are better off allowing widows and orphans (by changing the parameter
values) and manually correcting them in one way or another.}\cite[p.~250]{mittelbach2018forlorn}
Memoir's manual calls 10000 with a ragged bottom \enquote{the generally suggested method}, which \enquote{may well
result in some odd looking pages, especially if \apcs{raggedbottom} is not used}, and ends: \enquote{If all else
fails, the solution is either to live with the odd line or to reword the text.}\cite[§3.5, pp.~52--54]{wilson2025memoir}
For 1000, memoir's source gives the only stated reason: \enquote{The default values are 150 each, but we want
stronger discouragement.}\cite{wilson2025memoirsrc}
It gives no evidence, and no critique or defence of it was found.
Others: Osborne suggests 300 if a publisher prohibits such lines;\cite{osborne2010pmakeup} the
\texttt{widows-and-orphans} package offers 5000 for \enquote{avoid} and 10000 for \enquote{prevent}, and warns about
each case rather than deciding.\cite{mittelbach2023wao}
\texttt{lua-widow-control} lengthens a paragraph instead of leaving a gap.\cite{chernoff2022lwc}

\paragraph{Results: first line versus last line.}
The sources treat the two lines differently, while all three values are symmetric.
\emph{The Chicago Manual of Style} restricts only the last line, and only when short: a page should not
\enquote{begin with the last line of a paragraph unless it is full measure}, while \enquote{A page can, however,
end with the first line of a new paragraph}.\cite[§2.121]{chicago2024manual}
The 16th and 17th editions say the same.\cite[§2.113]{chicago2010manual}\cite[§2.116]{chicago2017manual}
The Press's sample house style agrees: the orphan \enquote{may fall at the foot of a page}, and the widow \enquote{may not fall at the top of a
page unless it is full measure}.\cite[fig.~A.3, item 7]{chicago2010manual}
The remedy is to run both facing pages a line long or short.\cite[§2.122]{chicago2024manual}
Wilson: \enquote{An orphan is not nearly so troubling to typographers as a widow.}\cite[§3.3.3, p.~52]{wilson2018notes}
Chernoff: orphans \enquote{tend not to decrease the legibility of a text as much as widows}, asserted without
data.\cite[p.~28]{chernoff2022lwc}
W3C: \enquote{many publishers now accept a single line of a paragraph before a page break}.\cite[§5.2]{cramer2016latinreq}
At second hand: Bringhurst, as quoted on a mailing list, says orphans \enquote{need not trouble the typographer}, but
that it is \enquote{the custom} to give a widow \enquote{one additional line for company}.\cite{izma2023widows}
Also second hand: Mittelbach quotes an older edition of the Manual as also forbidding a page to end with a first
line.\cite[p.~246]{mittelbach2018forlorn}
The U.S. Government Publishing Office permits both, the widow only \enquote{if absolutely necessary to maintain uniform
makeup and page depth} and the other more freely: \enquote{Paragraphs may start on the last line of a page whenever
necessary.}\cite[§§2.6--2.7]{gpo2016style}

\paragraph{Results: reader evidence.}
No study of widows or orphans in print was found, in either direction.
The closest is on screens: when sentences were split across screens, \enquote{neither variable significantly affects
comprehension}, but readers returned \enquote{to reread text significantly more}; this was read in abstract
only.\cite{dillon1990splitting}
A TUGboat review says that \enquote{modern experimental research questions the absolute importance} of typographic rules in
general.\cite[p.~285]{veytsman2011evidence}
A recent review warns that such \enquote{experiments rarely provide reliable evidence to support design
decisions}.\cite{vanderwaarde2025nineteen}
A blog reports of Wheildon's reader tests that \enquote{no readers said they were offended by -- or even aware of --
widows};\cite{harrison2020widows} this is second hand, the book was not read, and it concerns newspapers and magazines.
Wilson describes Wheildon's font studies as about 250 readers rating fonts.\cite[pp.~50--51]{wilson2018notes}
The confident claims about readers cite no data: \enquote{very distracting to the reader},\cite[p.~246]{mittelbach2018forlorn}
a widow \enquote{makes a paragraph harder to read},\cite[p.~28]{chernoff2022lwc} and single lines \enquote{should be
avoided}.\cite[Fig.~6]{schwen2020tensimple}
Butterick dissents from practice: \enquote{I only use it if widows and orphans are causing a visible
problem.}\cite{butterick2026widow}

\paragraph{Discussion and limitations.}
The mechanism test is one small document with one font and one page shape, and it tests when a line moves, not which
outcome is better.
The FAQ's statement is unmeasured and does not mention 1000.
Many sources argue from the look of the page, which a reader study would not necessarily settle.
Whether a class such as memoir embodies long typesetting experience is not something these documents show or refute.

\paragraph{Conclusion.}
The literature read does not support any particular value: no source compares 150, 1000 and 10000 for readers, and only one
gives a reason for 1000.
What is supported is the trade-off: 10000 removes such lines and can leave holes on inflexible pages, and any lower value leaves
them on such pages, with 150 and 1000 differing only when a page has a small amount of stretch.
The sources also treat the last line alone as worse than the first, at craft strength, and the Chicago Manual, in three editions and in its Press's
sample, permits the first line alone outright and the last line only at full measure.
None of these sources speaks of penalties, so they constrain the outcome and not a number.
What this package does with it: \cref{PageBreakDecisions}.

\paragraph{Future work.}
\emph{The \LaTeX\ Companion}, Bringhurst and Wheildon's \emph{Type \& Layout} should be read at source, as should
Chicago's earlier editions, to date the relaxation reported by Mittelbach.
Found in the databases but not obtained: Royster, \emph{Notes on Scholarly Journal Typesetting using Adobe InDesign} (2019); Hill,
\emph{The Creation of Typographic Specifications for Desktop Publishing Software} (2000); and Plass, \emph{Optimal Pagination
Techniques for Automatic Typesetting Systems} (1981), of which a scanned copy without text was found and not read.
The field lacks a controlled study of lone first and last lines in print or paged screens, which could show whether the asymmetry
reported by typographers has a reading effect; it would decide whether forbidding them buys anything for the space it costs.

\paragraph{Hit list.}
\Cref{tab:pb-hits} lists the database records; of those, Chernoff, Schwen, Veytsman and van der Waarde are cited here.

\section{Records bearing on the questions}\label{app:pb-hits}

^^A Author's note, short form: session heading-foot-of-page; regenerate with the
^^A command in the block at the top of this file, including the sed that removes
^^A the generated comment header.

\begingroup\footnotesize
\begin{longtable}{@{}%
  >{\raggedright\arraybackslash}p{0.44\textwidth}c%
  >{\raggedright\arraybackslash}p{0.13\textwidth}%
  >{\raggedright\arraybackslash}p{0.26\textwidth}@{}}
\caption{Records bearing on the claim, Page breaks around headings, lists, listings and paragraph lines (7 cited, 32 supporting and 2 qualifying or contradicting it, out of 2000 unique hits; the 818 adjacent and 1141 false hits are not listed). The last column gives the reason for each inclusion; for machine-classified rows the model's confidence is given. \enquote{Supports} means relevant to the questions by title or abstract, and an uncited row was not read. Unique hits count one work once across databases, so they are fewer than the records named in the text. Databases: OpenAlex (OA; not open access), IEEE Xplore (IEEE), Web of Science (WoS), Scopus, arXiv, HAL.}\label{tab:pb-hits}\\
\toprule
Title & Year & Provider & Reason \\
\midrule
\endfirsthead
\multicolumn{4}{@{}l}{\emph{\tablename~\thetable{} (cont.)}}\\
\toprule
Title & Year & Provider & Reason \\
\midrule
\endhead
\midrule \multicolumn{4}{r@{}}{\emph{continued on next page}}\\
\endfoot
\bottomrule
\endlastfoot
Attributes Influencing the Reading and Comprehension of Source Code -- Discussing Contradictory Evidence & 2018 & OA & \emph{cited}: reader evidence \\
Automatically removing widows and orphans with lua-widow-control & 2022 & OA & \emph{cited}: space at the page foot \\
Display Strategies for Program Browsing: Concepts and Experiment & 1986 & IEEE & \emph{cited}: reader evidence \\
Nineteen Questions to Evaluate Typographic Research: Chaff and Wheat & 2025 & OA & \emph{cited}: reader evidence \\
Ten simple rules for typographically appealing scientific texts & 2020 & OA & \emph{cited}: lines after a heading \\
Towards evidence-based typography: Literature review and experiment design & 2011 & OA & \emph{cited}: reader evidence \\
Typographic style is more than cosmetic & 1990 & OA & \emph{cited}: reader evidence \\
A general framework for globally optimized pagination & 2019 & Scopus & supports the claim (0.99) (title-or-abstract-only) \\
A technique for making structured programs more readable & 1978 & OA & supports the claim (0.72) (title-or-abstract-only) \\
A total-fit page-breaking algorithm with user-defined adjustment strategies & 2006 & Scopus & supports the claim (0.99) (title-or-abstract-only) \\
Algorithms and Listings & 2011 & OA & supports the claim (0.76) (title-or-abstract-only) \\
Breaking up & 1992 & OA & supports the claim (0.95) (title-or-abstract-only) \\
Definitive XSL-FO & 2003 & OA & supports the claim (0.99) (title-or-abstract-only) \\
Design principles for the enhanced presentation of computer program source text & 1986 & Scopus & supports the claim (0.86) (title-or-abstract-only) \\
Document layout optimization with automated paraphrasing & 2015 & Scopus & supports the claim (0.68) (title-or-abstract-only) \\
Extensible Stylesheet Language Requirements, Version 2.0, Working Draft 26 March 2008 & 2010 & OA & supports the claim (0.84) (title-or-abstract-only) \\
Formatting & 2010 & OA & supports the claim (0.97) (title-or-abstract-only) \\
Globally optimal page breaking with column balancing -- A case study & 2019 & Scopus & supports the claim (0.99) (title-or-abstract-only) \\
High-Quality Computer Typesetting for Text, Formulas, and Listings & 1969 & OA & supports the claim (0.78) (title-or-abstract-only) \\
Knižní úprava pro uživatele TeXu. Část druhá: Praxe & 1995 & OA & supports the claim (0.70) (title-or-abstract-only) \\
L A T E X tips for the academic & 2016 & OA & supports the claim (0.86) (title-or-abstract-only) \\
Line \& Page Breaks & 1995 & OA & supports the claim (0.99) (title-or-abstract-only) \\
Notes on Scholarly Journal Typesetting using Adobe InDesign & 2019 & OA & supports the claim (0.82) (title-or-abstract-only) \\
On the formatting of Pascal programs & 1977 & OA & supports the claim (0.77) (title-or-abstract-only) \\
Optimal pagination techniques for automatic typesetting systems & 1981 & OA & supports the claim (1.00) (title-or-abstract-only) \\
Ortotipografia: Sobre práticas e convenções tipográficas para o contexto editorial & 2022 & OA & supports the claim (0.76) (title-or-abstract-only) \\
Output routines extending \textbackslash{}plainoutput & 2025 & OA & supports the claim (0.76) (title-or-abstract-only) \\
Page Makeup by Postprocessing Text Formatter Output & 1989 & OA & supports the claim (0.99) (title-or-abstract-only) \\
PaperFit: Vision-in-the-Loop Typesetting Optimization for Scientific Documents & 2026 & OA & supports the claim (0.99) (title-or-abstract-only) \\
Sazba zdrojových kódů & 2012 & OA & supports the claim (0.91) (title-or-abstract-only) \\
Style guide for an Open Access publishing workflow for academic books & 2020 & OA & supports the claim (0.68) (title-or-abstract-only) \\
The Creation of typographic specifications for desktop publishing software & 2000 & OA & supports the claim (0.72) (title-or-abstract-only) \\
The Mac is not a typewriter: A style manual for creating professional-level type on your Macintosh & 1989 & OA & supports the claim (0.99) (title-or-abstract-only) \\
The Page Breaking Algorithm & 1993 & OA & supports the claim (0.99) (title-or-abstract-only) \\
Type Rules!: The Designer's Guide to Professional Typography & 2001 & OA & supports the claim (0.78) (title-or-abstract-only) \\
Typographic Style is More & 1990 & OA & supports the claim (0.99) (duplicate record of a cited source) \\
Updates to ``Automatically removing widows and orphans with lua-widow-control'', TUGboat 43:1 & 2022 & OA & supports the claim (0.99) (title-or-abstract-only) \\
Visible Languages for Program Visualization & 1986 & OA & supports the claim (0.98) (title-or-abstract-only) \\
Widows, Orphans and Bottom Balancing & 2012 & OA & supports the claim (1.00) (title-or-abstract-only) \\
A paradigm for programming style research & 1988 & OA & qualifies or contradicts the claim (0.95) (title-or-abstract-only) \\
Improving Code Reading and Comprehension on Large Displays & 2011 & OA & qualifies or contradicts the claim (0.83) (title-or-abstract-only) \\
\end{longtable}
\endgroup

